% answers.tex —— 1953 年·全国卷 · 参考答案与详细解析片段
% 由 tools/generate_projects.py 依来源总账自 1953/1953_数学真题_试题.md + 1953/1953_数学真题_解析正文.md + 1953/1953_数学真题_答案要点.md 生成，勿手改（改源 MD 后重跑）
\gkpart{一、小题}

\begin{answer}
\gksolhead{1}{}

\textbf{题目}：

（甲）解 $\dfrac{x+1}{x-1}+\dfrac{x-1}{x+1}=\dfrac{10}{3}$．

\Needspace{4\baselineskip}
\textbf{解析}

\Needspace{5\baselineskip}
通分并整理：

\nopagebreak[4]
\[
\frac{x+1}{x-1}+\frac{x-1}{x+1}=\frac{(x+1)^2+(x-1)^2}{x^2-1}=\frac{2x^2+2}{x^2-1}.
\]

\Needspace{5\baselineskip}
于是

\nopagebreak[4]
\[
\frac{2x^2+2}{x^2-1}=\frac{10}{3}
\implies 3(2x^2+2)=10(x^2-1)
\implies 6x^2+6=10x^2-10,
\]

即 $4x^2=16$，$x=\pm2$。

\gkanswer{答案：$x=\pm2$}

\end{answer}
\begin{answer}
\gksolhead{2}{}

\textbf{题目}：

（乙）若 $3x^2+kx+12=0$ 之二根相等，求 $k$．

\Needspace{4\baselineskip}
\textbf{解析}

\Needspace{5\baselineskip}
判别式

\nopagebreak[4]
\[
\Delta=k^2-4\cdot3\cdot12=k^2-144.
\]

二根相等当且仅当 $\Delta=0$，即 $k^2=144$，故 $k=\pm12$。

\gkanswer{答案：$k=\pm12$}

\end{answer}
\begin{answer}
\gksolhead{3}{}

\textbf{题目}：

（丙）求 $\begin{vmatrix}3&-1&1\\2&4&6\\7&0&5\end{vmatrix}$ 之值．

\Needspace{4\baselineskip}
\textbf{解析}

\Needspace{5\baselineskip}
按三阶行列式对角线法则展开：

\nopagebreak[4]
\[
\begin{vmatrix}3&-1&1\\2&4&6\\7&0&5\end{vmatrix}
=3\cdot4\cdot5+(-1)\cdot6\cdot7+1\cdot2\cdot0-1\cdot4\cdot7-(-1)\cdot2\cdot5-3\cdot6\cdot0
=60-42+0-28+10-0=0.
\]

\gkanswer{答案：$0$}

\end{answer}
\begin{answer}
\gksolhead{4}{}

\textbf{题目}：

（丁）求 $\log_{10}\dfrac{300}{7}+\log_{10}\dfrac{700}{3}+\log_{10}1$ 之值．

\Needspace{4\baselineskip}
\textbf{解析}

\Needspace{5\baselineskip}
由对数性质

\nopagebreak[4]
\[
\log_{10}\frac{300}{7}+\log_{10}\frac{700}{3}+\log_{10}1
=\log_{10}\left(\frac{300}{7}\cdot\frac{700}{3}\cdot1\right)
=\log_{10}10000=4.
\]

\gkanswer{答案：$4$}

\end{answer}
\begin{answer}
\gksolhead{5}{}

\textbf{题目}：

（戊）求 $\tan(870^\circ)$．

\Needspace{4\baselineskip}
\textbf{解析}

\nopagebreak[4]
\[
\tan870^\circ=\tan(870^\circ-2\times360^\circ)=\tan150^\circ
=\tan(180^\circ-30^\circ)=-\tan30^\circ=-\frac{\sqrt3}{3}.
\]

\gkanswer{答案：$-\dfrac{\sqrt3}{3}$}

\end{answer}
\begin{answer}
\gksolhead{6}{}

\textbf{题目}：

（己）若 $\cos2x=\dfrac12$，求 $x$ 之通值．

\Needspace{4\baselineskip}
\textbf{解析}

\nopagebreak[4]
\[
\cos2x=\frac12\implies 2x=2n\pi\pm\frac\pi3\implies x=n\pi\pm\frac\pi6\quad(n\in\mathbb Z).
\]

\gkanswer{答案：$n\pi\pm\dfrac\pi6$（$n$ 为整数）}

\end{answer}
\begin{answer}
\gksolhead{7}{}

\textbf{题目}：

（庚）两三角形相似之条件为何？（把你所知道的都写出来）

\Needspace{4\baselineskip}
\textbf{解析}

两三角形 $\triangle ABC$ 与 $\triangle A'B'C'$ 相似，下列条件之一成立即可：

（i）$\angle A=\angle A'$，$\angle B=\angle B'$（两角对应相等）；

（ii）$\angle A=\angle A'$，$AB:A'B'=AC:A'C'$（两边对应成比例且夹角相等）；

（iii）$AB:A'B'=AC:A'C'=BC:B'C'$（三边对应成比例）。

\gkanswer{答案：(i) 两角对应相等；(ii) 两边对应成比例且夹角相等；(iii) 三边对应成比例。}

\end{answer}
\begin{answer}
\gksolhead{8}{}

\textbf{题目}：

（辛）长方体之长、宽、高为 $12$ 寸、$3$ 寸、$4$ 寸，求对角线之长．

\Needspace{4\baselineskip}
\textbf{解析}

\Needspace{5\baselineskip}
长方体对角线长

\nopagebreak[4]
\[
d=\sqrt{12^2+3^2+4^2}=\sqrt{144+9+16}=\sqrt{169}=13.
\]

\gkanswer{答案：$13$ 寸}

\end{answer}
\begin{answer}
\gksolhead{9}{}

\textbf{题目}：

（壬）垂直三棱柱之高为 $6$ 寸，底面三边之长为 $3$ 寸、$4$ 寸、$5$ 寸，求体积．

\Needspace{4\baselineskip}
\textbf{解析}

\Needspace{5\baselineskip}
底面三边为 $3,4,5$，满足 $3^2+4^2=5^2$，故底面是直角三角形，面积

\nopagebreak[4]
\[
S=\frac12\cdot3\cdot4=6.
\]

棱柱体积 $=S\cdot h=6\times6=36$。

\gkanswer{答案：$36$ 立方寸}

\end{answer}
\begin{answer}
\gksolhead{10}{}

\textbf{题目}：

（癸）球之表面积为 $36\pi$ 方寸，求体积．

\Needspace{4\baselineskip}
\textbf{解析}

\Needspace{5\baselineskip}
设球半径为 $R$，则

\nopagebreak[4]
\[
4\pi R^2=36\pi\implies R^2=9\implies R=3.
\]

\Needspace{5\baselineskip}
体积

\nopagebreak[4]
\[
V=\frac43\pi R^3=\frac43\pi\cdot27=36\pi.
\]

\gkanswer{答案：$36\pi$ 立方寸}

\par\medskip

\end{answer}
\gkpart{二、大题}

\begin{answer}
\gksolhead{11}{}

\textbf{题目}：

（二）解 $\begin{cases}x^2-2xy+3y^2=9,\\4x^2-5xy+6y^2=30.\end{cases}$

\Needspace{4\baselineskip}
\textbf{解析}

\Needspace{5\baselineskip}
两式消去常数项：第一式乘 $10$ 减第二式乘 $3$，得

\nopagebreak[4]
\[
10(x^2-2xy+3y^2)-3(4x^2-5xy+6y^2)=0,
\]

\nopagebreak[4]
\[
-2x^2-5xy+12y^2=0\quad\Longleftrightarrow\quad 2x^2+5xy-12y^2=0.
\]

\Needspace{5\baselineskip}
左边分解因式

\nopagebreak[4]
\[
(x+4y)(2x-3y)=0,
\]

故 $x+4y=0$ 或 $2x-3y=0$。

\Needspace{5\baselineskip}
当 $x=-4y$ 时，代入 $x^2-2xy+3y^2=9$：

\nopagebreak[4]
\[
16y^2+8y^2+3y^2=9\implies 27y^2=9\implies y^2=\frac13,
\]

$y=\pm\dfrac{\sqrt3}{3}$，对应 $x=\mp\dfrac{4\sqrt3}{3}$，得两解 $\left(\dfrac{4\sqrt3}{3},-\dfrac{\sqrt3}{3}\right)$、$\left(-\dfrac{4\sqrt3}{3},\dfrac{\sqrt3}{3}\right)$。

\Needspace{5\baselineskip}
当 $2x=3y$，即 $x=\dfrac{3y}{2}$ 时，代入：

\nopagebreak[4]
\[
\frac{9y^2}{4}-3y^2+3y^2=9\implies \frac{9y^2}{4}=9\implies y^2=4,
\]

$y=\pm2$，对应 $x=\pm3$，得两解 $(3,2)$、$(-3,-2)$。

\gkanswer{答案：$\left(\dfrac{4\sqrt3}{3},-\dfrac{\sqrt3}{3}\right)$，$\left(-\dfrac{4\sqrt3}{3},\dfrac{\sqrt3}{3}\right)$，$(3,2)$，$(-3,-2)$}

\end{answer}
\begin{answer}
\gksolhead{12}{}

\textbf{题目}：

（三）设题目：

  (1) 化简 $\sqrt{\dfrac{12}{25}}+\sqrt[4]{90000}+\sqrt[6]{\dfrac{64}{27}}$．

  (2) 求 $\left(2x^3+\dfrac1x\right)^{12}$ 之展开式中之常数项．

\Needspace{4\baselineskip}
\textbf{解析}

\gksubpart{(1)}

\Needspace{5\baselineskip}
逐项化简：

\nopagebreak[4]
\[
\sqrt{\frac{12}{25}}=\frac{2\sqrt3}{5},\qquad
\sqrt[4]{90000}=\sqrt[4]{3^2\cdot10^4}=10\sqrt3,\qquad
\sqrt[6]{\frac{64}{27}}=\sqrt[6]{\frac{2^6}{3^3}}=\frac{2}{\sqrt3}=\frac{2\sqrt3}{3}.
\]

\Needspace{5\baselineskip}
三项相加：

\nopagebreak[4]
\[
\frac{2\sqrt3}{5}+10\sqrt3+\frac{2\sqrt3}{3}
=\left(\frac25+10+\frac23\right)\sqrt3
=\frac{6+150+10}{15}\sqrt3
=\frac{166\sqrt3}{15}.
\]

\gksubpart{(2)}

\Needspace{5\baselineskip}
二项展开通项

\nopagebreak[4]
\[
T_{r+1}=C_{12}^r(2x^3)^{12-r}\left(\frac1x\right)^r
=C_{12}^r2^{12-r}x^{36-3r-r}
=C_{12}^r2^{12-r}x^{36-4r}.
\]

\Needspace{5\baselineskip}
常数项要求 $36-4r=0$，即 $r=9$，故

\nopagebreak[4]
\[
T_{10}=C_{12}^9\cdot2^3=C_{12}^3\cdot8=220\cdot8=1760.
\]

\gkanswer{答案（1）：$\dfrac{166\sqrt3}{15}$}

\gkanswer{答案（2）：$1760$}

\end{answer}
\begin{answer}
\gksolhead{13}{}

\textbf{题目}：

（四）锐角三角形 $ABC$ 之三高线为 $AD$，$BE$，$CF$，垂心为 $H$；求证 $HD$ 平分 $\angle EDF$．

\Needspace{4\baselineskip}
\textbf{解析}

\Needspace{5\baselineskip}
由 $AD\perp BC$、$BE\perp CA$，得 $\angle ADB=\angle AEB=90^\circ$，故 $A,B,D,E$ 四点共圆，于是

\nopagebreak[4]
\[
\angle ADE=\angle ABE.
\]

同理，$F,B,C,E$ 四点共圆，得 $\angle FBE=\angle FCE$；$C,A,F,D$ 四点共圆，得 $\angle FCA=\angle FDA$。

\Needspace{5\baselineskip}
又 $BE$ 上有点 $B,E,H$，$CF$ 上有点 $C,F,H$，且 $\angle ABE=\angle FBE$、$\angle FCE=\angle FCA$（同角）。综上

\nopagebreak[4]
\[
\angle ADE=\angle ABE=\angle FBE=\angle FCE=\angle FCA=\angle FDA.
\]

故 $\angle ADE=\angle FDA$，即 $AD$（亦即 $HD$ 所在直线）平分 $\angle EDF$。

\gkanswer{答案：见解析版}

\end{answer}
\begin{answer}
\gksolhead{14}{}

\textbf{题目}：

（五）已知三角形的两个角为 $45^\circ$ 及 $60^\circ$，而其夹边长 $1$ 尺；求最小边之长及面积．

\Needspace{4\baselineskip}
\textbf{解析}

\Needspace{5\baselineskip}
设两已知角为 $\angle B=45^\circ$、$\angle C=60^\circ$，夹边 $a=BC=1$ 尺。则

\nopagebreak[4]
\[
\angle A=180^\circ-45^\circ-60^\circ=75^\circ.
\]

\Needspace{5\baselineskip}
最小边对最小角 $\angle B=45^\circ$，即 $AC=b$。由正弦定理

\nopagebreak[4]
\[
\frac{b}{\sin45^\circ}=\frac{1}{\sin75^\circ}
\implies b=\frac{\sin45^\circ}{\sin75^\circ}.
\]

\Needspace{5\baselineskip}
而

\nopagebreak[4]
\[
\sin75^\circ=\sin(45^\circ+30^\circ)=\sin45^\circ\cos30^\circ+\cos45^\circ\sin30^\circ
=\frac{\sqrt6+\sqrt2}{4}.
\]

\Needspace{5\baselineskip}
于是

\nopagebreak[4]
\[
b=\frac{\sqrt2/2}{(\sqrt6+\sqrt2)/4}
=\frac{2\sqrt2}{\sqrt6+\sqrt2}
=\frac{2\sqrt2(\sqrt6-\sqrt2)}{6-2}
=\frac{\sqrt{12}-2}{2}
=\frac{2\sqrt3-2}{2}
=\sqrt3-1.
\]

\Needspace{5\baselineskip}
面积

\nopagebreak[4]
\[
S=\frac12\cdot a\cdot b\sin C
=\frac12\cdot1\cdot(\sqrt3-1)\cdot\sin60^\circ
=\frac12(\sqrt3-1)\cdot\frac{\sqrt3}{2}
=\frac{3-\sqrt3}{4}.
\]

\gkanswer{答案：最小边 $\sqrt3-1$ 尺，面积 $\dfrac{3-\sqrt3}{4}$ 平方尺}

\end{answer}